% Kosten- und Finanzierungsplan für Förderanträge.
% Positionen aus den Arbeitspaketen der Projektbeschreibung, docs/COMPLIANCE.md
% und docs/funding-outline.md ("Numbers to collect"). Alle Beträge sind Platzhalter.
% Förderobergrenzen aus .orchestration/runs/2026-10-23-dossier/scout-applications.md;
% archivierte (A), sekundäre (S) und zu prüfende (TV) Angaben stehen nur mit \todo{}.
% Bauen aus docs/dossier/antrag/:
%   latexmk -pdf -interaction=nonstopmode -outdir=../build finanzplan.tex
\documentclass[11pt,a4paper,parskip=half-]{scrartcl}
% Shared preamble for the Anspruch dossier and its companion documents.
% Load after \documentclass with \input{preamble} (from docs/dossier/)
% or % Shared preamble for the Anspruch dossier and its companion documents.
% Load after \documentclass with \input{preamble} (from docs/dossier/)
% or % Shared preamble for the Anspruch dossier and its companion documents.
% Load after \documentclass with \input{preamble} (from docs/dossier/)
% or \input{../preamble} (from docs/dossier/antrag/ or docs/dossier/partner/).
% Each document sets its own \documentclass and its own pdftitle via
% \hypersetup{pdftitle={...}}. Only packages that ship with MiKTeX.

\usepackage[T1]{fontenc}
\usepackage[utf8]{inputenc}
\usepackage{lmodern}
\usepackage[english,ngerman]{babel} % last option = main language: ngerman
\usepackage[a4paper,margin=25mm]{geometry}
\usepackage{booktabs}
\usepackage{longtable}
\usepackage{array}
\usepackage{enumitem}
\usepackage{xcolor}
\usepackage{microtype}
\usepackage[autostyle=true,german=quotes]{csquotes}
\usepackage{marginnote}
\usepackage[hidelinks,unicode,pdflang={de-DE}]{hyperref}

% ---------------------------------------------------------------------------
% Colours
\definecolor{blankbg}{RGB}{255,236,153}   % highlighted placeholder
\definecolor{todofg}{RGB}{170,40,40}      % "to confirm" margin note

% \blank{Text}: a highlighted placeholder "[Text]" for values the sender
% fills in before sending. Usable in running text and in table cells.
\newcommand{\blank}[1]{\begingroup\setlength{\fboxsep}{1pt}%
  \colorbox{blankbg}{\strut[#1]}\endgroup}

% \todo{Text}: a "zu bestätigen" note in the margin. Works in running text,
% lists and longtable cells (via marginnote, not \marginpar).
% Defined with \providecommand so that it does not clash if a document
% loads todonotes on its own (todonotes then wins).
\providecommand{\todo}[1]{\marginnote{\raggedright\scriptsize
  \textcolor{todofg}{\textbf{Zu bestätigen:} #1}}}
\setlength{\marginparwidth}{20mm}

% \draftmark: header text for every page of a draft.
\newcommand{\draftmark}{Entwurf, 23.10.2026}

% Draft header for scrartcl / scrreprt (KOMA) and article / report.
% Letters (scrlttr2, letter) place \draftmark themselves.
\makeatletter
\@ifclassloaded{scrartcl}{\def\anspruch@koma{}}{}
\@ifclassloaded{scrreprt}{\def\anspruch@koma{}}{}
\@ifclassloaded{article}{\def\anspruch@std{}}{}
\@ifclassloaded{report}{\def\anspruch@std{}}{}
\ifdefined\anspruch@koma
  \usepackage[headsepline]{scrlayer-scrpage}
  \clearpairofpagestyles
  \ohead{\small\draftmark}
  \ihead{\small\headmark}
  \cfoot*{\pagemark}
  \automark[section]{section}
  \pagestyle{scrheadings}
\fi
\ifdefined\anspruch@std
  \usepackage{fancyhdr}
  \pagestyle{fancy}
  \fancyhf{}
  \fancyhead[R]{\small\draftmark}
  \fancyhead[L]{\small\nouppercase{\leftmark}}
  \fancyfoot[C]{\thepage}
  \fancypagestyle{plain}{\fancyhf{}\fancyhead[R]{\small\draftmark}%
    \fancyfoot[C]{\thepage}\renewcommand{\headrulewidth}{0pt}}
\fi
\makeatother

% Lists: compact, as in a German Antrag.
\setlist{itemsep=0.2em,topsep=0.4em}

% Tables: a bit more air, ragged-right paragraph columns.
\renewcommand{\arraystretch}{1.15}
\newcolumntype{L}[1]{>{\raggedright\arraybackslash}p{#1}}
 (from docs/dossier/antrag/ or docs/dossier/partner/).
% Each document sets its own \documentclass and its own pdftitle via
% \hypersetup{pdftitle={...}}. Only packages that ship with MiKTeX.

\usepackage[T1]{fontenc}
\usepackage[utf8]{inputenc}
\usepackage{lmodern}
\usepackage[english,ngerman]{babel} % last option = main language: ngerman
\usepackage[a4paper,margin=25mm]{geometry}
\usepackage{booktabs}
\usepackage{longtable}
\usepackage{array}
\usepackage{enumitem}
\usepackage{xcolor}
\usepackage{microtype}
\usepackage[autostyle=true,german=quotes]{csquotes}
\usepackage{marginnote}
\usepackage[hidelinks,unicode,pdflang={de-DE}]{hyperref}

% ---------------------------------------------------------------------------
% Colours
\definecolor{blankbg}{RGB}{255,236,153}   % highlighted placeholder
\definecolor{todofg}{RGB}{170,40,40}      % "to confirm" margin note

% \blank{Text}: a highlighted placeholder "[Text]" for values the sender
% fills in before sending. Usable in running text and in table cells.
\newcommand{\blank}[1]{\begingroup\setlength{\fboxsep}{1pt}%
  \colorbox{blankbg}{\strut[#1]}\endgroup}

% \todo{Text}: a "zu bestätigen" note in the margin. Works in running text,
% lists and longtable cells (via marginnote, not \marginpar).
% Defined with \providecommand so that it does not clash if a document
% loads todonotes on its own (todonotes then wins).
\providecommand{\todo}[1]{\marginnote{\raggedright\scriptsize
  \textcolor{todofg}{\textbf{Zu bestätigen:} #1}}}
\setlength{\marginparwidth}{20mm}

% \draftmark: header text for every page of a draft.
\newcommand{\draftmark}{Entwurf, 23.10.2026}

% Draft header for scrartcl / scrreprt (KOMA) and article / report.
% Letters (scrlttr2, letter) place \draftmark themselves.
\makeatletter
\@ifclassloaded{scrartcl}{\def\anspruch@koma{}}{}
\@ifclassloaded{scrreprt}{\def\anspruch@koma{}}{}
\@ifclassloaded{article}{\def\anspruch@std{}}{}
\@ifclassloaded{report}{\def\anspruch@std{}}{}
\ifdefined\anspruch@koma
  \usepackage[headsepline]{scrlayer-scrpage}
  \clearpairofpagestyles
  \ohead{\small\draftmark}
  \ihead{\small\headmark}
  \cfoot*{\pagemark}
  \automark[section]{section}
  \pagestyle{scrheadings}
\fi
\ifdefined\anspruch@std
  \usepackage{fancyhdr}
  \pagestyle{fancy}
  \fancyhf{}
  \fancyhead[R]{\small\draftmark}
  \fancyhead[L]{\small\nouppercase{\leftmark}}
  \fancyfoot[C]{\thepage}
  \fancypagestyle{plain}{\fancyhf{}\fancyhead[R]{\small\draftmark}%
    \fancyfoot[C]{\thepage}\renewcommand{\headrulewidth}{0pt}}
\fi
\makeatother

% Lists: compact, as in a German Antrag.
\setlist{itemsep=0.2em,topsep=0.4em}

% Tables: a bit more air, ragged-right paragraph columns.
\renewcommand{\arraystretch}{1.15}
\newcolumntype{L}[1]{>{\raggedright\arraybackslash}p{#1}}
 (from docs/dossier/antrag/ or docs/dossier/partner/).
% Each document sets its own \documentclass and its own pdftitle via
% \hypersetup{pdftitle={...}}. Only packages that ship with MiKTeX.

\usepackage[T1]{fontenc}
\usepackage[utf8]{inputenc}
\usepackage{lmodern}
\usepackage[english,ngerman]{babel} % last option = main language: ngerman
\usepackage[a4paper,margin=25mm]{geometry}
\usepackage{booktabs}
\usepackage{longtable}
\usepackage{array}
\usepackage{enumitem}
\usepackage{xcolor}
\usepackage{microtype}
\usepackage[autostyle=true,german=quotes]{csquotes}
\usepackage{marginnote}
\usepackage[hidelinks,unicode,pdflang={de-DE}]{hyperref}

% ---------------------------------------------------------------------------
% Colours
\definecolor{blankbg}{RGB}{255,236,153}   % highlighted placeholder
\definecolor{todofg}{RGB}{170,40,40}      % "to confirm" margin note

% \blank{Text}: a highlighted placeholder "[Text]" for values the sender
% fills in before sending. Usable in running text and in table cells.
\newcommand{\blank}[1]{\begingroup\setlength{\fboxsep}{1pt}%
  \colorbox{blankbg}{\strut[#1]}\endgroup}

% \todo{Text}: a "zu bestätigen" note in the margin. Works in running text,
% lists and longtable cells (via marginnote, not \marginpar).
% Defined with \providecommand so that it does not clash if a document
% loads todonotes on its own (todonotes then wins).
\providecommand{\todo}[1]{\marginnote{\raggedright\scriptsize
  \textcolor{todofg}{\textbf{Zu bestätigen:} #1}}}
\setlength{\marginparwidth}{20mm}

% \draftmark: header text for every page of a draft.
\newcommand{\draftmark}{Entwurf, 23.10.2026}

% Draft header for scrartcl / scrreprt (KOMA) and article / report.
% Letters (scrlttr2, letter) place \draftmark themselves.
\makeatletter
\@ifclassloaded{scrartcl}{\def\anspruch@koma{}}{}
\@ifclassloaded{scrreprt}{\def\anspruch@koma{}}{}
\@ifclassloaded{article}{\def\anspruch@std{}}{}
\@ifclassloaded{report}{\def\anspruch@std{}}{}
\ifdefined\anspruch@koma
  \usepackage[headsepline]{scrlayer-scrpage}
  \clearpairofpagestyles
  \ohead{\small\draftmark}
  \ihead{\small\headmark}
  \cfoot*{\pagemark}
  \automark[section]{section}
  \pagestyle{scrheadings}
\fi
\ifdefined\anspruch@std
  \usepackage{fancyhdr}
  \pagestyle{fancy}
  \fancyhf{}
  \fancyhead[R]{\small\draftmark}
  \fancyhead[L]{\small\nouppercase{\leftmark}}
  \fancyfoot[C]{\thepage}
  \fancypagestyle{plain}{\fancyhf{}\fancyhead[R]{\small\draftmark}%
    \fancyfoot[C]{\thepage}\renewcommand{\headrulewidth}{0pt}}
\fi
\makeatother

% Lists: compact, as in a German Antrag.
\setlist{itemsep=0.2em,topsep=0.4em}

% Tables: a bit more air, ragged-right paragraph columns.
\renewcommand{\arraystretch}{1.15}
\newcolumntype{L}[1]{>{\raggedright\arraybackslash}p{#1}}

\hypersetup{pdftitle={Anspruch: Kosten- und Finanzierungsplan}}

\begin{document}
{\LARGE\sffamily\bfseries Anspruch: Kosten- und Finanzierungsplan\par}
\blank{Name der Antragstellerin oder des Antragstellers} \textperiodcentered{}
Programm: \blank{Programm und Spur} \textperiodcentered{}
Laufzeit: \blank{Förderbeginn} bis \blank{Förderende}\par

\section*{Vor dem Versand ausfüllen}
\begin{quote}\raggedright
\blank{Name der Antragstellerin oder des Antragstellers} \textperiodcentered{}
\blank{Programm und Spur} \textperiodcentered{}
\blank{Förderbeginn} \textperiodcentered{} \blank{Förderende} \textperiodcentered{}
\blank{Monate} und \blank{Betrag} in jeder Zeile \textperiodcentered{}
\blank{Vercel-Tarif} \textperiodcentered{} \blank{Plausible-Tarif} \textperiodcentered{}
\blank{Summe} \textperiodcentered{}
\blank{Beantragte Förderung} \textperiodcentered{} \blank{Eigenanteil} \textperiodcentered{}
\blank{Partneranteil} \textperiodcentered{} \blank{Partner}
\end{quote}
Alle Beträge in EUR. Die Positionen folgen den Arbeitspaketen AP1 bis AP5 der Projektbeschreibung.
Welche Kostenarten ein Programm anerkennt, ist vor dem Ausfüllen beim Programm zu klären.

\section*{Kostenplan}
\begin{longtable}{@{}L{0.25\linewidth}L{0.41\linewidth}L{0.14\linewidth}L{0.12\linewidth}@{}}
\toprule
\textbf{Position} & \textbf{Beschreibung} & \textbf{Monate / Einheit} & \textbf{Betrag} \\
\midrule
\endhead
\bottomrule
\endlastfoot
1 Personal & Projektleitung, Entwicklung, Koordination der Arbeitspakete (Gründerin oder Gründer) &
  \blank{Monate} & \blank{Betrag} \\
2 Muttersprachliche Prüfung (AP1) & Je eine Person für Türkisch, Russisch, Polnisch und Arabisch,
  die das deutsche Pflegesystem kennt &
  4 Sprachen, \blank{Monate} & \blank{Betrag} \\
3 Rechtliche Prüfung (AP2) & Impressum und Datenschutzerklärung; Auftragsverarbeitungsverträge mit
  Vercel und Plausible; Prüfung zur Orientierungshilfe Digitale Dienste der DSK; Lizenz und Doppellizenz &
  \blank{Monate} & \blank{Betrag} \\
4 Hosting und Zählung & Vercel, Tarif \blank{Vercel-Tarif}; Plausible, Tarif \blank{Plausible-Tarif}
  (Eigenschaften je Partner brauchen nach Anbieterangabe den Business-Tarif)%
  \todo{Tarife und Preise bei Anmeldung prüfen; Vercel Hobby nur nicht kommerziell} &
  \blank{Monate} & \blank{Betrag} \\
5 Anpassung PNOG (AP4) & Nur falls das Gesetz beschlossen wird: Regelwerk, Quellen, Tests, Texte, PDF &
  \blank{Monate} & \blank{Betrag} \\
6 Pilotbegleitung (AP3) & Einrichtung der Partnerlinks, monatliche Rückmeldegespräche, Material für Beratende &
  \blank{Monate} & \blank{Betrag} \\
7 Auswertung (AP5) & Monatliche Summenzahlen je Partner, Strichlisten, Abschlussbericht &
  \blank{Monate} & \blank{Betrag} \\
8 Verwaltung & Buchhaltung, Berichte an den Fördergeber &
  \blank{Monate} & \blank{Betrag} \\
\midrule
\textbf{Summe} & & & \blank{Summe} \\
\end{longtable}

\section*{Finanzierungsplan}
\begin{longtable}{@{}L{0.36\linewidth}L{0.30\linewidth}L{0.26\linewidth}@{}}
\toprule
\textbf{Quelle} & \textbf{Anmerkung} & \textbf{Betrag} \\
\midrule
\endhead
\bottomrule
\endlastfoot
Beantragte Förderung & \blank{Programm und Spur} & \blank{Beantragte Förderung} \\
Eigenanteil & & \blank{Eigenanteil} \\
Anteil von Partnern & \blank{Partner} & \blank{Partneranteil} \\
\midrule
\textbf{Summe} (gleich Summe Kostenplan) & & \blank{Summe} \\
\end{longtable}

\section*{Obergrenzen der Programme}
Nur Angaben, die der Scout vom 23.10.2026 an der Quelle geprüft hat, stehen ohne Randnotiz.
\begin{itemize}
  \item \textbf{Prototype Fund:} \enquote{max.\ 950\,h für eine Person oder max.\ 1.900\,h für Teams.
        Die Maximalförderung beträgt 47.500 für eine Person oder 95.000 für Teams}, über sechs Monate.
        Laut Flyer \enquote{bis zu 95.000 in 6 Monaten oder 158.000 in 10 Monaten} mit Second Stage.
        Die Obergrenze der Second Stage selbst ist nicht belegt.%
        \todo{Archiv J02; für J03 bestätigen}
  \item \textbf{§ 123 SGB XI, Bund:} bis zu 30 Mio.\ EUR im Jahr aus dem Ausgleichsfonds, hälftig
        kofinanziert durch Land oder Kommune, 2025 bis 2028; höchstens vier Jahre, Ende spätestens
        31.12.2029.
  \item \textbf{§ 123 SGB XI, Niedersachsen:} höchstens 1 Mio.\ EUR je Vorhaben; Landesanteil
        45\,\% für Kommunen, 50\,\% für andere; die Pflegeversicherung gibt denselben Betrag dazu.
  \item \textbf{§ 123 SGB XI, Hessen:} mindestens 30.000 EUR im Jahr; Personalkosten höchstens nach
        TV-H oder TVöD; Investitionen nicht förderfähig; gefördert bis 31.12.2028.
  \item \textbf{§ 123 SGB XI, Berlin:} 200.000 EUR im Jahr für alle Vorhaben zusammen; die Frist
        für 2026 ist vorbei.
  \item \textbf{Deutsche Fernsehlotterie:} Gesamtkosten mindestens 50.000 EUR.
  \item \textbf{Aktion Mensch:} Eigenanteil je nach Programm.\todo{Sekundärquelle (10\,\% Beratung,
        20\,\% Investitionen); aktuelles Merkblatt prüfen}
\end{itemize}

\end{document}
